\section{Model and post-failure estimation bias}\label{sec:model}
We introduce the standard IBLT and the ideal peeling model, define the conditional miss rate required by the two-round protocol, characterize failure at low load, and explain how the standard IBLT's post-failure bias depends on sign composition. Appendix~\ref{app:notations} collects the notation used throughout the paper.

\subsection{The standard IBLT and the ideal peeling model}
The two sides hold sets $A$ and $B$ and seek their symmetric difference $A\triangle B$, whose size $d=|A\triangle B|$ is unknown in advance. The standard IBLT consists of $M$ cells, each with three fields: count, keySum, and checksum. To insert a key $x$, the cell hash selects $k$ distinct cells from $[M]$ and adds $1$, $x$, and $\mathrm{chk}(x)$ to their respective fields. Each side builds a table, and the tables are subtracted cell by cell. Shared keys contribute identical values to the same cells in both tables, so their contributions cancel in all three fields. Only the contributions of the $d$ differing keys remain. A key present only in $A$ has direction $s_i=+1$, and a key present only in $B$ has direction $s_i=-1$. We denote the numbers of keys in the two directions by $d_+$ and $d_-$. The ordered pair $(d_+,d_-)$ is the \emph{sign composition}, and $b=d_+-d_-=|A|-|B|$ is the difference between the set sizes. We use two extreme sign compositions throughout: a \emph{one-sided difference}, where all differing keys come from one side ($|b|=d$), and a \emph{balanced difference}, where both sides have the same number of differing keys ($d_+=d_-$, $b=0$). Both correspond to common synchronization scenarios. A one-sided difference arises in one-way synchronization when only one side inserts or deletes entries. A balanced difference arises in two-way synchronization when each side modifies a small number of entries and the insertions and deletions roughly cancel. Practical loads usually lie between these extremes, depending on how much each side changes. Decoding proceeds by peeling. We repeatedly find a pure cell containing one unrecovered differing key, recover the key by multiplying keySum by count, verify it with the checksum, and remove its contribution from all $k$ cells it selects. Removal may make other cells pure. Peeling continues until no pure cell remains.

Under ideal purity (a cell with count $\pm1$ contains exactly one unrecovered differing key), peeling success depends entirely on the mapping of differing keys to cells. View the $M$ cells as vertices and the $k$ cells selected by each differing key as a hyperedge. The difference table then corresponds to a random $k$-uniform hypergraph: the $d$ differing keys independently and uniformly choose $k$-subsets $S_i\subseteq[M]$, with different keys allowed to choose the same subset. We denote the unsigned mapping with key labels by $\mathcal G$ and the overlap between two edges by $H_{ij}=|S_i\cap S_j|$, the number of cells shared by differing keys $i$ and $j$. On $\mathcal G$, pure cells are exactly the unsigned degree-one vertices. Peeling repeatedly deletes edges incident to these vertices. The deletion order does not affect the terminal state, which is always the hypergraph's 2-core. A nonempty 2-core causes decoding failure, denoted by event $F$, with failure probability $\Pr(F)$. Thus $F$ depends only on the unsigned mapping $\mathcal G$. When the load $d/M$ is below the critical value $\alpha_c$ (about $0.818$ for $k=3$ and $0.772$ for $k=4$\cite{molloy,dw}), peeling succeeds with high probability. Rare failures arise from small cores such as duplicate edge pairs, and their probability decays polynomially in $M$. In the critical window of width of order $M^{-1/2}$ around $\alpha_c$, the failure probability rises from near zero to near one\cite{dm}. Above $\alpha_c$, peeling fails almost surely.

\subsection{The two-round protocol and the conditional miss rate}
Because the capacity $M$ must be fixed before $d$ is known, the two-round protocol\cite{wu2026selfsizing} first sends a small table of fixed size $M_1$. For a small difference, peeling succeeds in the first round. After a first-round failure, the receiver obtains an upper bound $U$ on $d$ from the difference table's count array and chooses the second-round size $m(U)$ based on $U$. The second round re-encodes and decodes using a fresh, independent cell hash.

Let $C_v$ denote the count of cell $v$ in the difference table. Following \cite{wu2026selfsizing}, we define the estimator $\widehat d$ of $d$ using the count sum $\sum_vC_v$ and sum of squares $\sum_vC_v^2$ saved before peeling begins:
\begin{equation}
\widehat d=\frac{1}{k(1-k/M)}\biggl(\sum_vC_v^2-\frac1M\Bigl(\sum_vC_v\Bigr)^2\biggr),\label{eq:dhat}
\end{equation}
and denote its unconditional standard deviation scale by $s(n)$:
\begin{equation}
s(n)=\sqrt{\frac{2n(n-1)}{M-1}}.\label{eq:s}
\end{equation}
For the standard IBLT, $\widehat d$ is exactly unbiased for $d$ and has variance $s(d)^2$ for every sign composition. This expectation is unconditional: it averages over all runs, regardless of failure $F$. If we evaluated the upper bound $U$ over all runs, adding a few standard deviations to $\widehat d$ would control the underestimate probability. In the two-round protocol, a successful first round ends the protocol immediately, without computing $U$. Only after decoding failure (event $F$) does the receiver compute $U$ to size the second-round table. If $U<d$, the second-round capacity falls below the prescribed margin, risking another failure. We call the event $U<d$ a \emph{miss} and its probability conditioned on first-round failure, $\Pr(U<d\mid F)$, the \emph{conditional miss rate}. The sizing rule must therefore satisfy
\begin{equation}
\Pr(U<d\mid F)\le u,\label{eq:target}
\end{equation}
where $u$ is the target risk. In deployment, we can choose $0.01$ or $0.05$, corresponding to a 1 \% or 5 \% underestimate risk.

The two-round protocol incurs a miss only when the first round fails and the upper bound underestimates $d$. This event has \emph{joint miss rate} $\Pr(F,\,U<d)=\Pr(F)\Pr(U<d\mid F)$. Controlling the conditional miss rate at $u$ bounds the joint miss rate by $u\Pr(F)$, both in overload, where failure is almost certain, and at low load, where it is rare. For failure samples without a miss ($U\ge d$), the second-round table has sufficient capacity margin to decode with high probability. Section~\ref{sec:protocol} gives the complete decomposition of the total two-round failure rate.

\subsection{Failure at low load and the planted distribution}\label{sec:plant}
\eqref{eq:target} must hold at all loads. Low load requires separate treatment. When $d$ is much smaller than $M$, the failure probability decays polynomially in $M$. Failure is then rare, making the distribution conditioned on $F$ difficult to analyze directly and requiring many simulated mappings for each failure sample. The structure of failure in this region is simple. Two edges with the same mapping form a duplicate edge pair. A single edge cannot form a 2-core, and two edges form a 2-core exactly when they are a duplicate edge pair. When $d=o(M)$, 2-cores with three or more edges are negligible relative to duplicate edge pairs. Failure is therefore approximately equivalent to the existence of a duplicate edge pair. Conditioning on $F$ is close to planting a duplicate edge pair in a random mapping. We thus define the \emph{planted distribution}: choose a pair of keys uniformly, assign both the same uniformly random $k$-subset, and let all remaining keys choose their mappings independently and uniformly.

\begin{proposition}[Failure is determined by duplicate edge pairs when $d=o(M)$]\label{prop:plant}
Fix $k\ge3$ and let $2\le d=o(M)$. The failure probability is asymptotically equal to the expected number of duplicate edge pairs: $\Pr(F)\sim\binom d2/\binom Mk$. The total variation distance between the distribution of $\mathcal G$ conditioned on $F$ and the planted distribution tends to zero.
\end{proposition}

Goodrich and Mitzenmacher\cite{gm} identify duplicate edge pairs as the main source of IBLT failure. The finite-length analysis of Yugawa and Wadayama\cite{yw} likewise shows that stopping sets of size 2 dominate. Proposition~\ref{prop:plant} extends this observation to a total variation approximation of the entire mapping for all $d=o(M)$. Lemma~\ref{lem:plant} in Appendix~\ref{app:graph} gives the quantitative error bound and proof. Proposition~\ref{prop:plant} supports both theory and experiments. It replaces conditioning on failure with the planted distribution, under which direct calculations are straightforward. The limit theorems in Sections~\ref{sec:measurement} and \ref{sec:inversion} use this approximation to handle conditioning on failure when $d=o(M)$. The risk decomposition in Section~\ref{sec:protocol} also uses it to estimate the first-round failure probability. For experiments, it enables direct sampling from the planted distribution, where every sample yields a failure. Section~\ref{sec:exp-setup} uses this planted sampling to evaluate low-load configurations beyond the reach of direct sampling and checks its finite-size accuracy where both methods are feasible.

\subsection{Post-failure bias of the standard IBLT}
Controlling \eqref{eq:target} is difficult because the failure event $F$ selects the mapping $\mathcal G$. Failure samples concentrate on mappings with large overlap $H_{ij}$, where differing keys share many cells. At low load, these mappings mainly contain duplicate edge pairs (Proposition~\ref{prop:plant}). In the critical window and at higher load, they contain large cores with dense overlap. Any statistic that depends on $H_{ij}$ therefore deviates from its unconditional distribution when conditioned on $F$.

In the standard IBLT, the difference-table count is $C_v=\sum_{i:v\in S_i}s_i$, and each differing key's coding coefficient is its fixed set direction $s_i\in\{+1,-1\}$. In the sum of $C_v^2$, the diagonal terms contribute exactly $kd$. The estimation error comes entirely from the cross terms $s_is_jH_{ij}$ between keys sharing cells. Excess overlap raises $\widehat d$ for same-sign pairs and lowers $\widehat d$ for opposite-sign pairs. The conditional relative bias of $\widehat d$ thus factors exactly into two terms.
\begin{proposition}[Bias decomposition of the standard IBLT]\label{prop:projection}
In the standard IBLT, for any sign composition $(d_+,d_-)$,
\begin{equation}
\frac{\E[\widehat d\mid F]}{d}-1=c_F\Bigl(\frac{b^2}{d}-1\Bigr),\qquad c_F=\frac{\E[H_{12}\mid F]-k^2/M}{k(1-k/M)}.\label{eq:projection}
\end{equation}
\end{proposition}
The proof is in Appendix~\ref{app:std}. Since $F$ is invariant under permutations of key labels, $\E[H_{ij}\mid F]$ is identical for all $i\ne j$. The formula uses differing keys $1$ and $2$ as a representative pair. The factor $c_F$ depends only on failure selection of the mapping and measures the excess of inter-edge overlap in the failure sample over its unconditional mean $k^2/M$. The factor $b^2/d-1$ depends only on the sign composition $(d_+,d_-)$. For a one-sided difference ($|b|=d$), every intersecting edge pair has sign product $+1$, so $b^2/d-1=d-1$. All excess overlap contributes with the same sign, fully amplifying its effect on $\widehat d$. For a balanced difference ($b=0$), there are $d/2$ fewer same-sign pairs than opposite-sign pairs, giving a relative bias of $-c_F$, whose magnitude is $1/(d-1)$ of the one-sided value.
\begin{figure}[!htbp]
\centering\includegraphics[width=\linewidth]{../figure/fig-bias.pdf}
\caption{Post-failure relative bias of the standard IBLT versus $(b/d)^2$ ($M=16384$). (a) $k=3$. (b) $k=6$. Each color denotes a failure probability. All configurations lie in the critical window: $d/M$ is 0.812--0.822 for $k=3$ and 0.633--0.639 for $k=6$. Points are measurements, and lines follow \eqref{eq:projection}.}
\label{fig:bias}
\end{figure}

Figure~\ref{fig:bias} shows this decomposition at $M=16384$ for five failure probabilities (0.10 to 0.90) each for $k=3$ and $k=6$. Within each configuration, all five sign compositions share the same batch of mappings and therefore identical peeling outcomes. Since $b^2/d-1=d(b/d)^2-1$, the relative bias at each failure probability is linear in $(b/d)^2$, with slope determined by $c_F$. We determine $c_F$ from the one-sided point. The remaining predictions of \eqref{eq:projection} differ from the measurements by at most twice the standard error of the difference. For $k=3$ at failure probability 0.10, the one-sided relative bias is $+0.70\%$, corresponding to a bias of about 0.64 standard deviations $s(d)$. As the failure probability increases, the failure sample approaches the full sample, weakening selection and reducing $|c_F|$. The sign of $c_F$ also depends on $k$. For $k=6$, $c_F<0$, so the one-sided failure sample is biased low, with relative bias $-0.43\%$ at failure probability 0.10 (Figure~\ref{fig:bias}(b)). The bias changes sign with $k$ and varies with sign composition according to \eqref{eq:projection}. A fixed correction therefore cannot offset it (Figure~\ref{fig:composition}).

In practice, the receiver does not know $(d_+,d_-)$ in advance. The difference table's count array alone cannot reveal whether a larger observed $\widehat d$ reflects a larger true $d$ or failure-induced bias for a one-sided difference. The upper-bound rule in \cite{wu2026selfsizing} therefore requires simulation calibration by load interval outside deep overload. By \eqref{eq:projection}, the bias arises because the standard IBLT uses the set direction $s_i$ as the accumulation sign in count, with sign composition determining the factor that amplifies $c_F$. For a balanced difference, this factor is $-1$, leaving a relative bias of $-c_F$. This suggests encoding differing keys with independent, uniform random signs so that every cross term has zero mean for each mapping, while preserving peeling's ability to recover each key's original set direction. The post-failure estimator would then be unbiased for every sign composition. We develop this idea in the random-sign IBLT in Section~\ref{sec:design}.
